\textbf{Krátké shrnutí a kdy začít DPIA.} Před pilotním nasazením zkontrolujte datový tok; u pilotů a při sdílení studentských dat s externími službami doporučena DPIA/DPA \cite{kb_ai_101_tipu_pdf_04e4a115_page_8_1}.

\textbf{Rychlý pracovní postup pro pedagoga.}
\begin{enumerate}
\item Omezte sběr na nezbytné údaje; neukládejte ID, pokud není nutné \cite{kb_ai_101_tipu_pdf_04e4a115_page_8_1}.
\item Před sdílením odstraňte PII, anonymizujte/pseudonymizujte a vymažte citlivá metadata \cite{kb_ai_101_tipu_pdf_04e4a115_page_8_1}.
\item Šifrujte a omezte přístupy; pro piloty ověřte DPIA/DPA a zajistěte lidský dohled \cite{kb_ai_101_tipu_pdf_04e4a115_page_8_1,kb_malinka_chatgpt_ve_vyuce_dva_roky_pote_2024_pdf_90bc3aaf_page_12_1}.
\end{enumerate}

\textbf{Hlavní položky DPIA a mitigace.}
\begin{itemize}
\item Popis účelu, datový tok, úložiště a přístupy; kategorie údajů a hodnocení rizik; mitigace (minimalizace, pseudonymizace, šifrování, DPA) a plán dohledu.
\end{itemize}

\textbf{Příklady a tip.} „Kvíz v Moodle“: anonymní skóre, externí export jen s DPA; „Projekt s kódem“: odstraňovat metadata a pseudonymizovat \cite{kb_ai_101_tipu_pdf_04e4a115_page_8_1}. Tip: fotografie, citlivá data nebo automatizované rozhodování bez lidského dohledu indikují potřebu DPIA a konzultaci \cite{kb_ai_101_tipu_pdf_04e4a115_page_8_1,kb_malinka_chatgpt_ve_vyuce_dva_roky_pote_2024_pdf_90bc3aaf_page_12_1}.